% Panorama de MCA_TOP EXPANDIDO: el mismo de panorama_mca_top.tex, con los
% bloques chicos abiertos a nivel de operación, como en el OSC:
%   - el disparo y el asignador (flanco, abre, la ranura de menor índice);
%   - la línea de base (el IIR con su holdoff);
%   - la contabilidad (los relojes del teselado, la ocupación y los períodos
%     muertos, cada uno con lo que lo habilita).
% Las ranuras, los trapecios, los divisores, el zoom y los histogramas quedan
% como submódulos. El disparo omite el pendiente de la decimación (un flanco
% entre dos muestras válidas espera a la siguiente): está en
% panorama_mca_window_expandido.tex.
% Las redes con nombre se conectan por nombre: el trazo corto que entra a un pin
% lleva el nombre de la red.
% Fuente: modulos/mca/rtl/mca_top.sv, mca_window.sv y mca_baseline.sv,
% instanciado como en rtl/mine/tops/mca_red_pitaya_top.sv (3 ranuras, 3 formas).
% Estilo: docs/figuras/estilo_tikz.tex; los colores por función, los mismos de
% panorama_mca_top.tex.
% El epígrafe no va en la figura: está en panorama_mca_top_expandido.txt.
%
% Compilar:   pdflatex panorama_mca_top_expandido.tex
%             pdftocairo -svg panorama_mca_top_expandido.pdf panorama_mca_top_expandido.svg
%
\documentclass[border=0pt,tikz]{standalone}
\usepackage[utf8]{inputenc}
\usepackage[T1]{fontenc}
\usepackage{lmodern}
\usepackage[none]{hyphenat}
% Estilo común de las figuras TikZ (panoramas y diagramas de bloques).
%
% Cada figura lo incluye con % Estilo común de las figuras TikZ (panoramas y diagramas de bloques).
%
% Cada figura lo incluye con % Estilo común de las figuras TikZ (panoramas y diagramas de bloques).
%
% Cada figura lo incluye con \input{../../figuras/estilo_tikz.tex} (las figuras
% viven en docs/<módulo>/figuras/). Un solo lugar para la paleta, los estilos de
% bloque y el formato de los parámetros: si cambia algo acá, cambia en todas.
%
% La paleta es la misma de docs/figuras/estilo_svg.py y de PALETA en
% modulos/mk/datasheet.py; si se toca un color, se toca en los tres.
%
% Requiere: \usepackage{tikz} (lo da la clase standalone con la opción tikz).

\usetikzlibrary{arrows.meta,positioning,calc,fit,backgrounds,shapes.multipart,
                decorations.markings,shadows}

% ---------------------------------------------------------------- colores ---
\definecolor{azul}{RGB}{108,172,228}      % acondicionamiento
\definecolor{verde}{RGB}{110,214,142}     % captura / histogramas en BRAM
\definecolor{amarillo}{RGB}{252,232,160}  % trigger
\definecolor{naranja}{RGB}{252,206,160}   % camino a DDR
\definecolor{grisr}{RGB}{226,226,226}     % compartido / infraestructura
\definecolor{rojo}{RGB}{240,128,128}      % bus del sistema

% ---------------------------------------------------------------- estilos ---
\tikzset{
  >={Stealth[length=2.4mm,width=1.8mm]},
  % fondo blanco siempre: el SVG se ve igual en modo claro y oscuro
  fondo/.style   = {show background rectangle, inner frame sep=8pt,
                    background rectangle/.style={fill=white, draw=none}},
  bloque/.style  = {draw=black, line width=0.6pt, rounded corners=2pt,
                    align=center, text width=2.2cm, minimum height=1.15cm,
                    inner sep=2pt, font=\scriptsize, fill=#1},
  datos/.style   = {bloque=azul!35},
  trig/.style    = {bloque=amarillo},
  captura/.style = {bloque=verde!55},
  ddr/.style     = {bloque=naranja!70},
  comp/.style    = {bloque=grisr},
  bus/.style     = {bloque=rojo!45},
  % lógica escrita dentro del módulo, sin submódulo propio
  etapa/.style   = {bloque=#1, line width=0.5pt},
  etapa/.default = verde!20,
  % lo que el bitstream no trae (atado a cero, stub, salida muda)
  mudo/.style    = {bloque=white, draw=black!45, dashed, text=black!55},
  %
  % BLOQUE CON PARÁMETROS, en el formato del datasheet: el encabezado lleva el
  % nombre y los parámetros con que lo instancia el padre (NOMBRE=valor, en
  % monoespaciada, uno por línea); el cuerpo, más claro, lo que hace.
  %   \node[datos, param={azul!35}{azul!12}] {\blk{nombre}{} \prm{A=1\\B=2}
  %                                            \nodepart{two} función};
  param/.style 2 args = {rectangle split, rectangle split parts=2,
                         rectangle split part fill={#1, #2}},
  %
  flujo/.style   = {->, line width=0.8pt},
  ctrl/.style    = {->, line width=0.6pt, black!65},
  % la línea que pasa por encima en un cruce: un halo blanco la separa
  cruce/.style   = {preaction={draw=white, line width=3.4pt, -}},
  % marca de ancho de bus: barra oblicua a mitad de la flecha, con el número
  % encima. \draw[flujo, ancho=14] ...; para otra posición, ancho en=0.3.
  ancho/.style   = {decoration={markings, mark=at position \anchopos with {
                      \draw[line width=0.6pt, -] (-0.1,-0.14) -- (0.1,0.14);
                      \node[font=\scriptsize, above=1.5pt] {#1};}},
                    postaction={decorate}},
  % la misma marca en una flecha vertical: el número va a la derecha
  anchov/.style  = {decoration={markings, mark=at position \anchopos with {
                      \draw[line width=0.6pt, -] (-0.14,-0.1) -- (0.14,0.1);
                      \node[font=\scriptsize, right=3pt] {#1};}},
                    postaction={decorate}},
  ancho en/.code = {\def\anchopos{#1}},
  sen/.style     = {font=\scriptsize\ttfamily, inner sep=1.5pt},
  nodo/.style    = {font=\tiny\ttfamily, inner sep=1pt, black!75},
  pin/.style     = {font=\scriptsize, align=center, inner sep=1pt},
  grupo/.style   = {draw=black!55, dashed, line width=0.7pt, rounded corners=4pt},
  canal/.style   = {grupo},
  contenedor/.style = {draw=black!80, line width=1.0pt, rounded corners=6pt},
  % marca numerada: cada número es una fila de la tabla del documento
  marca/.style   = {circle, fill=black, text=white, font=\scriptsize\bfseries,
                    inner sep=0pt, minimum size=4.3mm},
  %
  % DECISIONES (docs/mca/vista_general_mca.md). Cada opción es un bloque, todas
  % entran a un mux (\muxdec, abajo) y el registro que elige entra por abajo.
  % la opción que trae el reset (el comportamiento histórico): borde doble
  defecto/.style = {double=white, double distance=0.8pt, line width=0.45pt},
  % rótulo de la decisión, el mismo Dn del panorama y del documento
  decision/.style = {circle, fill=black, text=white, font=\small\bfseries,
                     inner sep=0pt, minimum size=7mm},
  %
  % LÓGICA (docs/mca/figuras/contabilidad/ y los detalles de las decisiones)
  % registro: caja gris con el triángulo de reloj en el borde de abajo
  registro/.style = {draw=black, line width=0.6pt, fill=grisr, align=center,
                     font=\scriptsize, inner sep=2pt, minimum width=#1,
                     minimum height=0.85cm,
                     path picture={\draw[line width=0.5pt]
                       ([xshift=-1.1mm]path picture bounding box.south) --
                       ([yshift=1.6mm]path picture bounding box.south) --
                       ([xshift=1.1mm]path picture bounding box.south);}},
  registro/.default = 0.8cm,
  % condición de una cadena if / else-if: "sí" sale a la derecha, "no" abajo
  pregunta/.style = {draw=black, line width=0.5pt, rounded corners=7pt,
                     fill=white, font=\scriptsize\ttfamily, align=center,
                     inner sep=2.5pt, minimum height=0.62cm, text width=#1},
  pregunta/.default = 1.9cm,
  si/.style      = {font=\tiny, inner sep=1pt, above, pos=0.12},
  no/.style      = {font=\tiny, inner sep=1pt, left, pos=0.45},
  %
  titulo/.style  = {anchor=north west, font=\small, inner sep=4pt, align=left},
  rotulo/.style  = {anchor=north west, font=\scriptsize\itshape, black!70,
                    inner sep=2pt, align=left},
}

% submódulo: nombre en monoespaciada; lógica propia: nombre en romana
\newcommand{\blk}[2]{{\ttfamily\bfseries #1}\\[1pt]#2}
\newcommand{\etq}[2]{{\bfseries #1}\\[1pt]#2}
\newcommand{\punto}[1]{\fill (#1) circle (1.3pt);}

% Parámetros en el formato del datasheet: NOMBRE=valor en monoespaciada, una
% entrada por celda, hasta cuatro columnas. Una columna: \prm{A=1\\B=2};
% dos: \prm{A=1 & B=2\\C=3}.
\newcommand{\prm}[1]{{\ttfamily\fontsize{5.6}{6.6}\selectfont
  \begin{tabular}{@{}l@{\hspace{7pt}}l@{\hspace{7pt}}l@{\hspace{7pt}}l@{}}#1\end{tabular}}}

\def\anchopos{0.5}

% Mux de una decisión: trapecio con el lado ancho (las entradas) a la izquierda.
% #1 = nombre, #2 = centro, #3 = alto en cm. Deja tres coordenadas:
%   (#1-in)  centro del lado de entradas: las flechas llegan a (#1-in |- opcion.east)
%   (#1-out) la salida
%   (#1-sel) mitad del lado de abajo: por ahí entra el registro que elige
% El rótulo de adentro es opcional: \muxdec{m}{(4,0)}{3.6} y después
% \node[...] at (m) {...}.
\newcommand{\muxdec}[3]{%
  \path #2 coordinate (#1);
  \draw[line width=0.6pt, fill=grisr]
    ($(#1)+(-0.45,{#3/2})$) -- ($(#1)+(0.45,{#3/2-0.35})$) --
    ($(#1)+(0.45,{-#3/2+0.35})$) -- ($(#1)+(-0.45,{-#3/2})$) -- cycle;
  \coordinate (#1-in)  at ($(#1)+(-0.45,0)$);
  \coordinate (#1-out) at ($(#1)+(0.45,0)$);
  \coordinate (#1-sel) at ($(#1)+(0,{-#3/2+0.175})$);}

% El mismo mux rotado 90 grados, para un flujo de arriba hacia abajo: entradas
% por el lado ancho de arriba, salida abajo y el registro que elige por la
% izquierda. #1 = nombre, #2 = centro, #3 = ancho en cm. Deja (#1-in) centro del
% lado de entradas, (#1-out) la salida y (#1-sel) mitad del lado izquierdo.
\newcommand{\muxdecv}[3]{%
  \path #2 coordinate (#1);
  \draw[line width=0.6pt, fill=grisr]
    ($(#1)+({-#3/2},0.45)$) -- ($(#1)+({#3/2},0.45)$) --
    ($(#1)+({#3/2-0.35},-0.45)$) -- ($(#1)+({-#3/2+0.35},-0.45)$) -- cycle;
  \coordinate (#1-in)  at ($(#1)+(0,0.45)$);
  \coordinate (#1-out) at ($(#1)+(0,-0.45)$);
  \coordinate (#1-sel) at ($(#1)+({-#3/2+0.175},0)$);}

% Leyenda de colores + nota. #1 = posición, #2 = color/texto, #3 = nota.
\newcommand{\leyenda}[3]{%
  \begin{scope}[shift={(#1)}, every node/.style={font=\scriptsize}]
    \coordinate (t0) at (-0.5,0);
    \foreach \c/\t [count=\i] in {#2} {
      \pgfmathtruncatemacro\j{\i-1}
      \node[draw, fill=\c, minimum width=4mm, minimum height=3mm, inner sep=0pt,
            right=5mm of t\j.east, anchor=west] (l\i) {};
      \node[right=1.5mm of l\i] (t\i) {\t};
    }
    \node[anchor=north west, black!70, text width=19cm, align=left]
      at (-0.05,-0.3) {#3};
  \end{scope}}

% Lo que dicen los parámetros, para la nota de la leyenda
\newcommand{\notaparam}{\texttt{NOMBRE=valor}: parámetros con que el padre
  instancia el bloque (se omiten los que sólo pasan el del padre sin cambio).}
 (las figuras
% viven en docs/<módulo>/figuras/). Un solo lugar para la paleta, los estilos de
% bloque y el formato de los parámetros: si cambia algo acá, cambia en todas.
%
% La paleta es la misma de docs/figuras/estilo_svg.py y de PALETA en
% modulos/mk/datasheet.py; si se toca un color, se toca en los tres.
%
% Requiere: \usepackage{tikz} (lo da la clase standalone con la opción tikz).

\usetikzlibrary{arrows.meta,positioning,calc,fit,backgrounds,shapes.multipart,
                decorations.markings,shadows}

% ---------------------------------------------------------------- colores ---
\definecolor{azul}{RGB}{108,172,228}      % acondicionamiento
\definecolor{verde}{RGB}{110,214,142}     % captura / histogramas en BRAM
\definecolor{amarillo}{RGB}{252,232,160}  % trigger
\definecolor{naranja}{RGB}{252,206,160}   % camino a DDR
\definecolor{grisr}{RGB}{226,226,226}     % compartido / infraestructura
\definecolor{rojo}{RGB}{240,128,128}      % bus del sistema

% ---------------------------------------------------------------- estilos ---
\tikzset{
  >={Stealth[length=2.4mm,width=1.8mm]},
  % fondo blanco siempre: el SVG se ve igual en modo claro y oscuro
  fondo/.style   = {show background rectangle, inner frame sep=8pt,
                    background rectangle/.style={fill=white, draw=none}},
  bloque/.style  = {draw=black, line width=0.6pt, rounded corners=2pt,
                    align=center, text width=2.2cm, minimum height=1.15cm,
                    inner sep=2pt, font=\scriptsize, fill=#1},
  datos/.style   = {bloque=azul!35},
  trig/.style    = {bloque=amarillo},
  captura/.style = {bloque=verde!55},
  ddr/.style     = {bloque=naranja!70},
  comp/.style    = {bloque=grisr},
  bus/.style     = {bloque=rojo!45},
  % lógica escrita dentro del módulo, sin submódulo propio
  etapa/.style   = {bloque=#1, line width=0.5pt},
  etapa/.default = verde!20,
  % lo que el bitstream no trae (atado a cero, stub, salida muda)
  mudo/.style    = {bloque=white, draw=black!45, dashed, text=black!55},
  %
  % BLOQUE CON PARÁMETROS, en el formato del datasheet: el encabezado lleva el
  % nombre y los parámetros con que lo instancia el padre (NOMBRE=valor, en
  % monoespaciada, uno por línea); el cuerpo, más claro, lo que hace.
  %   \node[datos, param={azul!35}{azul!12}] {\blk{nombre}{} \prm{A=1\\B=2}
  %                                            \nodepart{two} función};
  param/.style 2 args = {rectangle split, rectangle split parts=2,
                         rectangle split part fill={#1, #2}},
  %
  flujo/.style   = {->, line width=0.8pt},
  ctrl/.style    = {->, line width=0.6pt, black!65},
  % la línea que pasa por encima en un cruce: un halo blanco la separa
  cruce/.style   = {preaction={draw=white, line width=3.4pt, -}},
  % marca de ancho de bus: barra oblicua a mitad de la flecha, con el número
  % encima. \draw[flujo, ancho=14] ...; para otra posición, ancho en=0.3.
  ancho/.style   = {decoration={markings, mark=at position \anchopos with {
                      \draw[line width=0.6pt, -] (-0.1,-0.14) -- (0.1,0.14);
                      \node[font=\scriptsize, above=1.5pt] {#1};}},
                    postaction={decorate}},
  % la misma marca en una flecha vertical: el número va a la derecha
  anchov/.style  = {decoration={markings, mark=at position \anchopos with {
                      \draw[line width=0.6pt, -] (-0.14,-0.1) -- (0.14,0.1);
                      \node[font=\scriptsize, right=3pt] {#1};}},
                    postaction={decorate}},
  ancho en/.code = {\def\anchopos{#1}},
  sen/.style     = {font=\scriptsize\ttfamily, inner sep=1.5pt},
  nodo/.style    = {font=\tiny\ttfamily, inner sep=1pt, black!75},
  pin/.style     = {font=\scriptsize, align=center, inner sep=1pt},
  grupo/.style   = {draw=black!55, dashed, line width=0.7pt, rounded corners=4pt},
  canal/.style   = {grupo},
  contenedor/.style = {draw=black!80, line width=1.0pt, rounded corners=6pt},
  % marca numerada: cada número es una fila de la tabla del documento
  marca/.style   = {circle, fill=black, text=white, font=\scriptsize\bfseries,
                    inner sep=0pt, minimum size=4.3mm},
  %
  % DECISIONES (docs/mca/vista_general_mca.md). Cada opción es un bloque, todas
  % entran a un mux (\muxdec, abajo) y el registro que elige entra por abajo.
  % la opción que trae el reset (el comportamiento histórico): borde doble
  defecto/.style = {double=white, double distance=0.8pt, line width=0.45pt},
  % rótulo de la decisión, el mismo Dn del panorama y del documento
  decision/.style = {circle, fill=black, text=white, font=\small\bfseries,
                     inner sep=0pt, minimum size=7mm},
  %
  % LÓGICA (docs/mca/figuras/contabilidad/ y los detalles de las decisiones)
  % registro: caja gris con el triángulo de reloj en el borde de abajo
  registro/.style = {draw=black, line width=0.6pt, fill=grisr, align=center,
                     font=\scriptsize, inner sep=2pt, minimum width=#1,
                     minimum height=0.85cm,
                     path picture={\draw[line width=0.5pt]
                       ([xshift=-1.1mm]path picture bounding box.south) --
                       ([yshift=1.6mm]path picture bounding box.south) --
                       ([xshift=1.1mm]path picture bounding box.south);}},
  registro/.default = 0.8cm,
  % condición de una cadena if / else-if: "sí" sale a la derecha, "no" abajo
  pregunta/.style = {draw=black, line width=0.5pt, rounded corners=7pt,
                     fill=white, font=\scriptsize\ttfamily, align=center,
                     inner sep=2.5pt, minimum height=0.62cm, text width=#1},
  pregunta/.default = 1.9cm,
  si/.style      = {font=\tiny, inner sep=1pt, above, pos=0.12},
  no/.style      = {font=\tiny, inner sep=1pt, left, pos=0.45},
  %
  titulo/.style  = {anchor=north west, font=\small, inner sep=4pt, align=left},
  rotulo/.style  = {anchor=north west, font=\scriptsize\itshape, black!70,
                    inner sep=2pt, align=left},
}

% submódulo: nombre en monoespaciada; lógica propia: nombre en romana
\newcommand{\blk}[2]{{\ttfamily\bfseries #1}\\[1pt]#2}
\newcommand{\etq}[2]{{\bfseries #1}\\[1pt]#2}
\newcommand{\punto}[1]{\fill (#1) circle (1.3pt);}

% Parámetros en el formato del datasheet: NOMBRE=valor en monoespaciada, una
% entrada por celda, hasta cuatro columnas. Una columna: \prm{A=1\\B=2};
% dos: \prm{A=1 & B=2\\C=3}.
\newcommand{\prm}[1]{{\ttfamily\fontsize{5.6}{6.6}\selectfont
  \begin{tabular}{@{}l@{\hspace{7pt}}l@{\hspace{7pt}}l@{\hspace{7pt}}l@{}}#1\end{tabular}}}

\def\anchopos{0.5}

% Mux de una decisión: trapecio con el lado ancho (las entradas) a la izquierda.
% #1 = nombre, #2 = centro, #3 = alto en cm. Deja tres coordenadas:
%   (#1-in)  centro del lado de entradas: las flechas llegan a (#1-in |- opcion.east)
%   (#1-out) la salida
%   (#1-sel) mitad del lado de abajo: por ahí entra el registro que elige
% El rótulo de adentro es opcional: \muxdec{m}{(4,0)}{3.6} y después
% \node[...] at (m) {...}.
\newcommand{\muxdec}[3]{%
  \path #2 coordinate (#1);
  \draw[line width=0.6pt, fill=grisr]
    ($(#1)+(-0.45,{#3/2})$) -- ($(#1)+(0.45,{#3/2-0.35})$) --
    ($(#1)+(0.45,{-#3/2+0.35})$) -- ($(#1)+(-0.45,{-#3/2})$) -- cycle;
  \coordinate (#1-in)  at ($(#1)+(-0.45,0)$);
  \coordinate (#1-out) at ($(#1)+(0.45,0)$);
  \coordinate (#1-sel) at ($(#1)+(0,{-#3/2+0.175})$);}

% El mismo mux rotado 90 grados, para un flujo de arriba hacia abajo: entradas
% por el lado ancho de arriba, salida abajo y el registro que elige por la
% izquierda. #1 = nombre, #2 = centro, #3 = ancho en cm. Deja (#1-in) centro del
% lado de entradas, (#1-out) la salida y (#1-sel) mitad del lado izquierdo.
\newcommand{\muxdecv}[3]{%
  \path #2 coordinate (#1);
  \draw[line width=0.6pt, fill=grisr]
    ($(#1)+({-#3/2},0.45)$) -- ($(#1)+({#3/2},0.45)$) --
    ($(#1)+({#3/2-0.35},-0.45)$) -- ($(#1)+({-#3/2+0.35},-0.45)$) -- cycle;
  \coordinate (#1-in)  at ($(#1)+(0,0.45)$);
  \coordinate (#1-out) at ($(#1)+(0,-0.45)$);
  \coordinate (#1-sel) at ($(#1)+({-#3/2+0.175},0)$);}

% Leyenda de colores + nota. #1 = posición, #2 = color/texto, #3 = nota.
\newcommand{\leyenda}[3]{%
  \begin{scope}[shift={(#1)}, every node/.style={font=\scriptsize}]
    \coordinate (t0) at (-0.5,0);
    \foreach \c/\t [count=\i] in {#2} {
      \pgfmathtruncatemacro\j{\i-1}
      \node[draw, fill=\c, minimum width=4mm, minimum height=3mm, inner sep=0pt,
            right=5mm of t\j.east, anchor=west] (l\i) {};
      \node[right=1.5mm of l\i] (t\i) {\t};
    }
    \node[anchor=north west, black!70, text width=19cm, align=left]
      at (-0.05,-0.3) {#3};
  \end{scope}}

% Lo que dicen los parámetros, para la nota de la leyenda
\newcommand{\notaparam}{\texttt{NOMBRE=valor}: parámetros con que el padre
  instancia el bloque (se omiten los que sólo pasan el del padre sin cambio).}
 (las figuras
% viven en docs/<módulo>/figuras/). Un solo lugar para la paleta, los estilos de
% bloque y el formato de los parámetros: si cambia algo acá, cambia en todas.
%
% La paleta es la misma de docs/figuras/estilo_svg.py y de PALETA en
% modulos/mk/datasheet.py; si se toca un color, se toca en los tres.
%
% Requiere: \usepackage{tikz} (lo da la clase standalone con la opción tikz).

\usetikzlibrary{arrows.meta,positioning,calc,fit,backgrounds,shapes.multipart,
                decorations.markings,shadows}

% ---------------------------------------------------------------- colores ---
\definecolor{azul}{RGB}{108,172,228}      % acondicionamiento
\definecolor{verde}{RGB}{110,214,142}     % captura / histogramas en BRAM
\definecolor{amarillo}{RGB}{252,232,160}  % trigger
\definecolor{naranja}{RGB}{252,206,160}   % camino a DDR
\definecolor{grisr}{RGB}{226,226,226}     % compartido / infraestructura
\definecolor{rojo}{RGB}{240,128,128}      % bus del sistema

% ---------------------------------------------------------------- estilos ---
\tikzset{
  >={Stealth[length=2.4mm,width=1.8mm]},
  % fondo blanco siempre: el SVG se ve igual en modo claro y oscuro
  fondo/.style   = {show background rectangle, inner frame sep=8pt,
                    background rectangle/.style={fill=white, draw=none}},
  bloque/.style  = {draw=black, line width=0.6pt, rounded corners=2pt,
                    align=center, text width=2.2cm, minimum height=1.15cm,
                    inner sep=2pt, font=\scriptsize, fill=#1},
  datos/.style   = {bloque=azul!35},
  trig/.style    = {bloque=amarillo},
  captura/.style = {bloque=verde!55},
  ddr/.style     = {bloque=naranja!70},
  comp/.style    = {bloque=grisr},
  bus/.style     = {bloque=rojo!45},
  % lógica escrita dentro del módulo, sin submódulo propio
  etapa/.style   = {bloque=#1, line width=0.5pt},
  etapa/.default = verde!20,
  % lo que el bitstream no trae (atado a cero, stub, salida muda)
  mudo/.style    = {bloque=white, draw=black!45, dashed, text=black!55},
  %
  % BLOQUE CON PARÁMETROS, en el formato del datasheet: el encabezado lleva el
  % nombre y los parámetros con que lo instancia el padre (NOMBRE=valor, en
  % monoespaciada, uno por línea); el cuerpo, más claro, lo que hace.
  %   \node[datos, param={azul!35}{azul!12}] {\blk{nombre}{} \prm{A=1\\B=2}
  %                                            \nodepart{two} función};
  param/.style 2 args = {rectangle split, rectangle split parts=2,
                         rectangle split part fill={#1, #2}},
  %
  flujo/.style   = {->, line width=0.8pt},
  ctrl/.style    = {->, line width=0.6pt, black!65},
  % la línea que pasa por encima en un cruce: un halo blanco la separa
  cruce/.style   = {preaction={draw=white, line width=3.4pt, -}},
  % marca de ancho de bus: barra oblicua a mitad de la flecha, con el número
  % encima. \draw[flujo, ancho=14] ...; para otra posición, ancho en=0.3.
  ancho/.style   = {decoration={markings, mark=at position \anchopos with {
                      \draw[line width=0.6pt, -] (-0.1,-0.14) -- (0.1,0.14);
                      \node[font=\scriptsize, above=1.5pt] {#1};}},
                    postaction={decorate}},
  % la misma marca en una flecha vertical: el número va a la derecha
  anchov/.style  = {decoration={markings, mark=at position \anchopos with {
                      \draw[line width=0.6pt, -] (-0.14,-0.1) -- (0.14,0.1);
                      \node[font=\scriptsize, right=3pt] {#1};}},
                    postaction={decorate}},
  ancho en/.code = {\def\anchopos{#1}},
  sen/.style     = {font=\scriptsize\ttfamily, inner sep=1.5pt},
  nodo/.style    = {font=\tiny\ttfamily, inner sep=1pt, black!75},
  pin/.style     = {font=\scriptsize, align=center, inner sep=1pt},
  grupo/.style   = {draw=black!55, dashed, line width=0.7pt, rounded corners=4pt},
  canal/.style   = {grupo},
  contenedor/.style = {draw=black!80, line width=1.0pt, rounded corners=6pt},
  % marca numerada: cada número es una fila de la tabla del documento
  marca/.style   = {circle, fill=black, text=white, font=\scriptsize\bfseries,
                    inner sep=0pt, minimum size=4.3mm},
  %
  % DECISIONES (docs/mca/vista_general_mca.md). Cada opción es un bloque, todas
  % entran a un mux (\muxdec, abajo) y el registro que elige entra por abajo.
  % la opción que trae el reset (el comportamiento histórico): borde doble
  defecto/.style = {double=white, double distance=0.8pt, line width=0.45pt},
  % rótulo de la decisión, el mismo Dn del panorama y del documento
  decision/.style = {circle, fill=black, text=white, font=\small\bfseries,
                     inner sep=0pt, minimum size=7mm},
  %
  % LÓGICA (docs/mca/figuras/contabilidad/ y los detalles de las decisiones)
  % registro: caja gris con el triángulo de reloj en el borde de abajo
  registro/.style = {draw=black, line width=0.6pt, fill=grisr, align=center,
                     font=\scriptsize, inner sep=2pt, minimum width=#1,
                     minimum height=0.85cm,
                     path picture={\draw[line width=0.5pt]
                       ([xshift=-1.1mm]path picture bounding box.south) --
                       ([yshift=1.6mm]path picture bounding box.south) --
                       ([xshift=1.1mm]path picture bounding box.south);}},
  registro/.default = 0.8cm,
  % condición de una cadena if / else-if: "sí" sale a la derecha, "no" abajo
  pregunta/.style = {draw=black, line width=0.5pt, rounded corners=7pt,
                     fill=white, font=\scriptsize\ttfamily, align=center,
                     inner sep=2.5pt, minimum height=0.62cm, text width=#1},
  pregunta/.default = 1.9cm,
  si/.style      = {font=\tiny, inner sep=1pt, above, pos=0.12},
  no/.style      = {font=\tiny, inner sep=1pt, left, pos=0.45},
  %
  titulo/.style  = {anchor=north west, font=\small, inner sep=4pt, align=left},
  rotulo/.style  = {anchor=north west, font=\scriptsize\itshape, black!70,
                    inner sep=2pt, align=left},
}

% submódulo: nombre en monoespaciada; lógica propia: nombre en romana
\newcommand{\blk}[2]{{\ttfamily\bfseries #1}\\[1pt]#2}
\newcommand{\etq}[2]{{\bfseries #1}\\[1pt]#2}
\newcommand{\punto}[1]{\fill (#1) circle (1.3pt);}

% Parámetros en el formato del datasheet: NOMBRE=valor en monoespaciada, una
% entrada por celda, hasta cuatro columnas. Una columna: \prm{A=1\\B=2};
% dos: \prm{A=1 & B=2\\C=3}.
\newcommand{\prm}[1]{{\ttfamily\fontsize{5.6}{6.6}\selectfont
  \begin{tabular}{@{}l@{\hspace{7pt}}l@{\hspace{7pt}}l@{\hspace{7pt}}l@{}}#1\end{tabular}}}

\def\anchopos{0.5}

% Mux de una decisión: trapecio con el lado ancho (las entradas) a la izquierda.
% #1 = nombre, #2 = centro, #3 = alto en cm. Deja tres coordenadas:
%   (#1-in)  centro del lado de entradas: las flechas llegan a (#1-in |- opcion.east)
%   (#1-out) la salida
%   (#1-sel) mitad del lado de abajo: por ahí entra el registro que elige
% El rótulo de adentro es opcional: \muxdec{m}{(4,0)}{3.6} y después
% \node[...] at (m) {...}.
\newcommand{\muxdec}[3]{%
  \path #2 coordinate (#1);
  \draw[line width=0.6pt, fill=grisr]
    ($(#1)+(-0.45,{#3/2})$) -- ($(#1)+(0.45,{#3/2-0.35})$) --
    ($(#1)+(0.45,{-#3/2+0.35})$) -- ($(#1)+(-0.45,{-#3/2})$) -- cycle;
  \coordinate (#1-in)  at ($(#1)+(-0.45,0)$);
  \coordinate (#1-out) at ($(#1)+(0.45,0)$);
  \coordinate (#1-sel) at ($(#1)+(0,{-#3/2+0.175})$);}

% El mismo mux rotado 90 grados, para un flujo de arriba hacia abajo: entradas
% por el lado ancho de arriba, salida abajo y el registro que elige por la
% izquierda. #1 = nombre, #2 = centro, #3 = ancho en cm. Deja (#1-in) centro del
% lado de entradas, (#1-out) la salida y (#1-sel) mitad del lado izquierdo.
\newcommand{\muxdecv}[3]{%
  \path #2 coordinate (#1);
  \draw[line width=0.6pt, fill=grisr]
    ($(#1)+({-#3/2},0.45)$) -- ($(#1)+({#3/2},0.45)$) --
    ($(#1)+({#3/2-0.35},-0.45)$) -- ($(#1)+({-#3/2+0.35},-0.45)$) -- cycle;
  \coordinate (#1-in)  at ($(#1)+(0,0.45)$);
  \coordinate (#1-out) at ($(#1)+(0,-0.45)$);
  \coordinate (#1-sel) at ($(#1)+({-#3/2+0.175},0)$);}

% Leyenda de colores + nota. #1 = posición, #2 = color/texto, #3 = nota.
\newcommand{\leyenda}[3]{%
  \begin{scope}[shift={(#1)}, every node/.style={font=\scriptsize}]
    \coordinate (t0) at (-0.5,0);
    \foreach \c/\t [count=\i] in {#2} {
      \pgfmathtruncatemacro\j{\i-1}
      \node[draw, fill=\c, minimum width=4mm, minimum height=3mm, inner sep=0pt,
            right=5mm of t\j.east, anchor=west] (l\i) {};
      \node[right=1.5mm of l\i] (t\i) {\t};
    }
    \node[anchor=north west, black!70, text width=19cm, align=left]
      at (-0.05,-0.3) {#3};
  \end{scope}}

% Lo que dicen los parámetros, para la nota de la leyenda
\newcommand{\notaparam}{\texttt{NOMBRE=valor}: parámetros con que el padre
  instancia el bloque (se omiten los que sólo pasan el del padre sin cambio).}

\usetikzlibrary{decorations.pathreplacing}
\tikzset{
  linea/.style   = {rectangle split, rectangle split parts=2},
  pre/.style     = {bloque=azul!35, param={azul!35}{azul!12}},
  est/.style     = {bloque=verde!55, param={verde!55}{verde!20}},
  pub/.style     = {bloque=naranja!70, param={naranja!70}{naranja!25}},
  publog/.style  = {etapa=naranja!30},
  varios/.style  = {double copy shadow={shadow xshift=2pt, shadow yshift=-2pt}},
  chico/.style   = {font=\tiny, inner sep=1pt, align=center},
  op/.style      = {oper, font=\scriptsize},
  caja/.style    = {desp, inner sep=2pt, minimum height=0.5cm},
  reg/.style     = {registro=#1, minimum height=0.55cm, font=\scriptsize\ttfamily},
  reg/.default   = 1.05cm,
  cnt/.style     = {registro=1.9cm, minimum height=0.48cm, font=\scriptsize\ttfamily},
  fondo f/.style = {draw=black!45, dashed, line width=0.5pt, rounded corners=3pt, fill=#1},
  rotf/.style    = {font=\scriptsize\itshape, black!75, inner sep=2pt},
  red/.style     = {font=\scriptsize\ttfamily, inner sep=1pt},
  y3/.style      = {logic gate input sep=0.2cm},
}
\newcommand{\red}[3][0.3]{\draw[flujo] ($(#2)+(-#1,0)$) node[left, red] {#3} -- (#2);}

\begin{document}
\begin{tikzpicture}[fondo]

\def\ytop{4.9}    % de donde bajan las entradas

% ============================================================ las regiones ===
\node[fondo f=grisr!45, fit={(6.75,-0.85) (15.2,3.65)}] (fcon) {};
\node[fondo f=amarillo!22, fit={(3.85,-1.95) (10.55,-4.5)}] (fasg) {};
\node[rotf, anchor=south west] at (fasg.south west) {disparo y asignador};
\node[rotf, anchor=north west] at (fcon.north west) {contabilidad};
\node[fondo f=azul!10, fit={(1.05,-5.55) (3.4,-11.0)}] (fbas) {};
\node[rotf, anchor=north west, align=left] at (fbas.north west) {\texttt{mca\_baseline}};
\node[fondo f=azul!10, fit={(4.45,-8.85) (9.45,-11.0)}] (fhol) {};
\node[rotf, anchor=south east] at (fhol.south east) {cuándo sigue la base};

% ======================================================== la contabilidad ===
% los relojes del teselado: cada ciclo con run suma en exactamente uno
\node[cnt] (crt) at (10.05,3.0) {REALTIME};
\node[cnt] (cve) at (10.05,2.4) {VETOTIME};
\node[cnt] (cwi) at (10.05,1.8) {T\_WIN};
\node[cnt] (cdi) at (10.05,1.2) {T\_DIV};
\node[cnt] (cra) at (10.05,0.6) {T\_REARM};
\node[cnt] (cli) at (10.05,0.0) {LIVETIME};
\red{crt.west}{run}
\red{cve.west}{veto}
\red{cwi.west}{$\neg$veto$\wedge$dt\_win}
\red{cdi.west}{$\neg$veto$\wedge$dt\_div}
\red{cra.west}{$\neg$veto$\wedge$puede$\wedge\neg$arm}
\red{cli.west}{$\neg$veto$\wedge$puede$\wedge$arm}
\draw[decorate, decoration={brace, amplitude=4pt}, line width=0.5pt]
  ($(cve.north east)+(0.08,0)$) -- ($(cli.south east)+(0.08,0)$);
\node[chico, anchor=north west, align=left] at (6.95,0 |- cli.south) {la llave: el teselado, un casillero por ciclo;\\ cada uno, $+1$ por ciclo, de 64 b};
% la ocupación, y lo que no entra en el teselado
\node[cnt] (coc) at (14.15,3.0) {T\_OCUP[$j$]};
\node[cnt] (cde) at (14.15,2.4) {DEADTIME};
\node[cnt] (ccl) at (14.15,1.8) {T\_CLEAR};
\red{coc.west}{$\neg$veto}
\red{cde.west}{$\neg$veto$\wedge$busy}
\red{ccl.west}{clear}
\node[chico, anchor=south east] at (coc.north east) {$j$ = \texttt{n\_ocup}};
% los períodos muertos: tramos sin ranura libre
\node[cnt] (cpe) at (14.15,1.0) {DEAD\_PERIODS};
\node[cnt] (cml) at (14.15,0.4) {MAX\_DEAD\_LEN};
\node[cnt] (cex) at (14.15,-0.2) {DEAD\_EXT};
\red{cpe.west}{$\uparrow$bloq}
\red{cml.west}{máx $L$}
\red{cex.west}{bloq$\wedge$arribo}
\node[chico, anchor=north east] at (15.15,0 |- cex.south) {bloq = $\neg$veto$\wedge$(dt\_win$\vee$dt\_div)};
% entradas de la contabilidad
\draw[flujo] (8.2,\ytop) node[above, sen] {veto\_i} -- (8.2,0 |- fcon.north);
\draw[flujo] (9.2,\ytop) node[above, sen] {trig\_arm\_i} -- (9.2,0 |- fcon.north);
\node[chico, anchor=north west, align=left] at (8.25,3.6) {$\to$ veto};
\node[chico, anchor=north west, align=left] at (9.25,3.6) {$\to$ arm};

% ================================================== mca_pulse_feature ===
% --- la muestra y la resta ---
\node[circle, draw, line width=0.6pt, fill=azul!15, minimum size=5.5mm, inner sep=0pt,
      font=\small] (res) at (4.2,-5.0) {$-$};
\node[registro, minimum width=0.85cm] (xq) at (5.3,-5.0) {\texttt{xc\_q}\\ \texttt{x\_q}};

% --- el disparo y el asignador, abiertos ---
\node[reg=0.95cm] (tq) at (4.95,-3.25) {trig\_q};
\node[puerta y=nin, y3, anchor=input 1] (edg) at (5.85,-2.25) {};
\draw[flujo] (4.15,\ytop) node[above, sen] {trig\_i} -- (4.15,-2.25) -- (edg.input 1);
\draw[flujo] (4.15,-3.25) -- (tq.west);
\punto{4.15,-3.25}
\draw[flujo] (tq.east) -| ($(edg.input 2)+(-0.2,0)$) -- (edg.input 2);
\red{edg.input 3}{val$\wedge$run}
\node[chico, anchor=south] at ($(edg.output)+(0.32,0)$) {disparo};
\node[puerta y=nin, y3, anchor=input 1] (abr) at (8.0,0 |- edg.output) {};
\draw[flujo] (edg.output) -- (abr.input 1);
\red{abr.input 2}{veto}
\red{abr.input 3}{puede}
\node[op] (pue) at (6.85,-3.4) {$\neq0$};
\red{pue.west}{libre}
\draw[flujo] (pue.east) -- ++(0.3,0) node[right, red] {puede};
\node[caja] (pri) at (7.0,-4.15) {menor\\ índice};
\red{pri.west}{libre}
\node[caja, anchor=west] (oh) at ($(pri.east)+(0.2,0)$) {$1 \ll$ sel};
\draw[flujo] (pri.east) -- (oh.west);
\node[puerta y, anchor=input 2] (aoh) at ($(oh.east)+(0.3,0)$) {};
\draw[flujo] (oh.east) -- (aoh.input 2);
\draw[flujo] (abr.output) -- ++(0.3,0) node[right, red] {abre};
\draw[flujo] ($(aoh.input 1)+(-0.2,0.45)$) node[above, red] {abre} |- (aoh.input 1);
\node[reg=0.85cm] (raq) at (10.25,0 |- aoh.output) {abre\_q};
\draw[flujo] (aoh.output) -- (raq.west);

% --- trapecios y ranuras ---
\node[pre, text width=1.85cm, varios] (trp) at (5.75,-6.45)
  {\blk{mca\_trap {\normalfont$\times$} 3}{}\nodepart{two}$k$, $m$ propios;\\ $M$ común};
\node[est, text width=3.0cm, varios] (acu) at (8.9,-6.05)
  {\blk{mca\_slot {\normalfont$\times$} 3}{}\nodepart{two}pico, $Q_\mathrm{tot}$, $Q_\mathrm{cola}$,\\
   largo, $t_\mathrm{pico}$, \texttt{dt\_next};\\ ventana de amplitud;\\ una meseta por forma};
\draw[flujo] (raq.south) -- node[chico, right, align=left] {+ marcas} (raq.south |- acu.north);
\node[grupo, fit={(3.6,-1.75) (10.7,-8.55)}] (gw) {};
\node[rotulo, anchor=north east] at (gw.north east) {\texttt{mca\_window}};

% --- mca_publish ---
\node[est, text width=2.8cm, varios] (div) at (12.45,-3.05)
  {\blk{mca\_div\_restore {\normalfont$\times$} 2}{}\nodepart{two}$Q_\mathrm{cola}/Q_\mathrm{tot} \to$ \texttt{F\_PSD}\\
   pico$/Q_\mathrm{tot} \to$ \texttt{F\_INVW}\\[1pt] {\tiny 16 b en 16 ciclos}};
\node[publog, text width=2.8cm, linea] (pubr) at (12.45,-6.25)
  {{\bfseries publicador}\nodepart{two}de a una cerrada,\\ en orden de cierre:\\ elige la forma,\\ arma el bus, libera\\[1pt] {\tiny 18 ciclos por evento}};
\node[grupo, fit={(10.95,-1.75) (14.1,-8.55)}] (gp) {};
\node[rotulo, anchor=south east] at (gp.south east) {\texttt{mca\_publish}};

% --- la línea de base, abierta: b_acc += dat - (b_acc >> k), base = b_acc >> k o fija ---
\node[op] (bme) at (2.3,-6.15) {$-$};
\node[op] (bma) at (2.3,-7.0) {$+$};
\node[reg=1.0cm] (bac) at (2.3,-7.9) {bl\_acc};
\node[chico, anchor=west, align=left] at (3.05,-7.65) {en:\\ sigue};
\node[caja] (bsh) at (2.3,-8.85) {$\gg k$};
\muxdecv{bmx}{(2.65,-9.75)}{1.1}
\node[reg=1.0cm] (bas) at (2.65,-10.55) {base};
\draw[flujo] (bme.south) -- (bma.north);
\draw[flujo] (bma.south) -- (bac.north);
\draw[flujo] (bac.south) -- (bsh.north);
\coordinate (jb) at ($(bac.south)!0.5!(bsh.north)$);
\draw[flujo] (jb) -- ++(0.65,0) |- (bma.east);
\punto{jb}
\draw[flujo] (bsh.west) -- ++(-0.6,0) |- (bme.west);
\draw[flujo] (bsh.south) -- (2.3,0 |- bmx-in);
\draw[flujo] (3.02,-9.05) node[above, red] {fija} -- (3.02,0 |- bmx-in);
\draw[ctrl] ($(bmx-sel)+(-0.4,0)$) node[left, red, black!80] {auto} -- (bmx-sel);
\draw[flujo] (bmx-out) -- (bas.north);
% --- cuándo sigue: ninguna ranura midiendo, y pasado el holdoff del último cierre ---
\node[reg=1.15cm] (bho) at (5.65,-9.65) {bl\_hold};
\node[chico, anchor=north, align=center] at (bho.south) {carga holdoff al cerrar;\\ $-1$ por muestra válida};
\node[op] (bce) at (6.9,-9.65) {$=0$};
\draw[flujo] (bho.east) -- (bce.west);
\red{bho.west}{cierra}
\node[puerta y=nnn, y3, anchor=input 3] (bsg) at (7.75,0 |- bce.east) {};
\draw[flujo] (bce.east) -- (bsg.input 3);
\draw[flujo] ($(bsg.input 1)+(-0.25,0.35)$) node[above, red] {val, bl\_track} |- (bsg.input 1);
\draw[flujo] ($(bsg.input 1)+(-0.25,0)$) |- (bsg.input 2);
\draw[flujo] (bsg.output) -- ++(0.25,0) node[right, red] {sigue};

\node[contenedor, draw=black!55, line width=0.7pt, fit={(0.9,-1.35) (14.25,-11.25)}] (pf) {};
\node[rotulo, anchor=north west] at (pf.north west)
  {\texttt{mca\_pulse\_feature}\; (solo cableado)};

% =============================================== del evento al histograma ===
\def\yc{-14.05}
\node[publog, text width=1.75cm, linea] (mux) at (1.9,\yc)
  {{\bfseries mux de eje}\nodepart{two}16:1 por eje,\\ \texttt{cfg\_sel\_*}};
\node[registro, minimum width=0.95cm] (fq) at (3.75,\yc) {\texttt{f\_q}\\ \texttt{inc}};
\node[pub, text width=1.8cm, varios] (zoo) at (5.75,\yc)
  {\blk{mca\_zoom {\normalfont$\times$} 3}{}\nodepart{two}feature $\to$ bin,\\ ventana $2^{-z}$};
\node[registro, minimum width=0.95cm] (pip) at (7.75,\yc) {\texttt{bin}\\ \texttt{inc}};
\node[pub, text width=2.75cm, varios] (hst) at (10.35,\yc)
  {\blk{mca\_hist {\normalfont$\times$} 2}{}\nodepart{two}1D: 8192 $\times$ 32 b\\
   2D: $128{\times}64 \times 32$ b};
\node[pub, text width=2.3cm] (dis) at (13.8,\yc)
  {\blk{mca\_discriminator}{}\nodepart{two}una feature\\ en $[\min, \max]$};

% ======================================================= registros y bus ===
\node[comp, text width=14.3cm, minimum height=0.95cm, linea] (reg) at (8.1,-16.9)
  {{\bfseries banco de registros}\nodepart{two}\texttt{cfg\_*} \;$\cdot$\; contadores de 64 b con sombra
   \;$\cdot$\; mux de lectura \;$\cdot$\; ack de latencia fija};
\node[bus, text width=1.8cm] (sb) at (8.1,-18.5) {\texttt{\bfseries sys\_bus}};

% ============================================================ caja mca_top
\node[contenedor, fit={(0.65,4.0) (15.55,-17.6)}] (top) {};
\node[titulo, anchor=south east, inner sep=2pt, align=left] at (top.north east)
  {\texttt{\bfseries mca\_top}\;\; {\scriptsize\itshape expandido, todo en \texttt{adc\_clk}}\\[1pt]
   \prm{DW=14 & H\_AW=13 & H2\_AW=7\\ PSD\_AW=6 & N\_ACC=3 & N\_FORMAS=3}};

% ============================================================== flechas ====
% --- la muestra: a la base y a la resta ---
\coordinate (jd) at (2.3,-5.0);
\node[sen, align=center, anchor=south] at (2.3,\ytop)
  {dat\_d\\[-1pt]\normalfont\scriptsize un canal, retrasado};
\draw[flujo, anchov=14, ancho en=0.15] (2.3,\ytop) -- (jd) -- (res.west);
\punto{jd}
\draw[flujo] (jd) -- (bme.north);
\draw[flujo] (bas.east) -- node[sen, above, pos=0.5] {base} (4.2,0 |- bas.east) -- (res.south);

% --- x = dat - base, registrada: saturada a las ranuras, con signo al banco ---
\draw[flujo] (res.east) -- (xq.west);
\draw[flujo, ancho=14, ancho en=0.4] (xq.east) -- node[sen, below, pos=0.4] {xc} (6.65,0 |- xq.east) |- (acu.west |- 0,-5.55);
\draw[flujo, anchov=15, ancho en=0.4] (xq.south) -- node[sen, left, pos=0.4] {x} (xq.south |- trp.north);
\draw[flujo] (trp.east) -- node[chico, below, fill=white] {3$\times$18} (trp.east -| acu.west);

% --- de las ranuras al publicador, y del publicador a los divisores ---
\draw[flujo] (acu.east |- 0,-5.35) -- node[chico, above, pos=0.5, fill=white] {cerradas} (pubr.west |- 0,-5.35);
\draw[ctrl] (pubr.west |- 0,-6.3) -- node[chico, below, pos=0.5, fill=white] {libera} (acu.east |- 0,-6.3);
\draw[flujo] (11.9,0 |- pubr.north) -- node[chico, left] {$Q$, pico} (11.9,0 |- div.south);
\draw[flujo] (13.0,0 |- div.south) -- node[chico, right] {cocientes} (13.0,0 |- pubr.north);

% --- de la ventana a la contabilidad ---
\draw[flujo] (10.4,0 |- gw.north) -- node[chico, right, align=left, pos=0.45]
  {dt\_win, dt\_div, puede,\\ busy, n\_ocup, arribo} (10.4,0 |- fcon.south);

% --- el bus de features: una barra, y lo que cuelga de ella ---
\def\ybus{-12.3}
\draw[line width=0.8pt] (pubr.south) -- (12.45,\ybus);
\draw[line width=2.2pt] (1.9,\ybus) -- (13.8,\ybus);
\node[sen, anchor=south, align=center] at (6.9,\ybus)
  {ev\_feat: 16 features $\times$ 16 b, uno por evento, registrado\\[-1pt]
   \normalfont\scriptsize + ev\_valid, ev\_psd\_ok};
\draw[flujo] (1.9,\ybus) -- (mux.north);
\draw[flujo] (13.8,\ybus) -- (dis.north);

% --- de la feature al bin y al histograma ---
\draw[flujo] (mux.east) -- (fq.west);
\draw[flujo] (fq.east) -- (zoo.west);
\draw[flujo] (zoo.east) -- (pip.west);
\draw[flujo] (pip.east) -- (hst.west);
\node[chico, anchor=north west] at (8.27,\yc-0.05) {13 b\\ 7+6 b};
\draw[flujo] (13.15,0 |- dis.north) -- (13.15,-12.95) -- node[sen, above, pos=0.62]
  {inc = ev\_valid $\wedge$ acepta} (3.75,-12.95) -- (fq.north);

% --- registros, contadores y configuración ---
\draw[flujo, cruce, anchov=32, ancho en=0.25] (hst.south) -- node[sen, right, pos=0.75] {rd\_data} (hst.south |- reg.north);
\draw[flujo] (fcon.east |- 0,1.4) -- (15.38,1.4) -- (15.38,0 |- reg.north);
\node[sen, anchor=south, rotate=90] at (15.36,0.2) {tiempos};
\draw[line width=0.8pt] (pf.east |- 0,-7.6) -- (15.38,-7.6);
\punto{15.38,-7.6}
\node[sen, anchor=south east] at (15.3,-7.55) {contadores};
\draw[ctrl, cruce] (12.15,0 |- reg.north) -- node[sen, black!80, right, pos=0.25] {cfg\_*} (12.15,0 |- pf.south);

% --- bus del sistema ---
\draw[<->, line width=0.8pt] (reg.south) -- (sb.north);
\draw[<->, line width=0.8pt] (sb.east) -- ++(1.2,0) node[right, sen, align=left]
  {M\_AXI\_GP0\\[-1pt]\normalfont\scriptsize$\leftrightarrow$ ARM};

% ================================================================= leyenda ===
\leyenda{0.75,-19.45}{azul!35/línea de base y conformado, amarillo!75/segmentación,
  verde!55/estadísticos del pulso}{}
\leyenda{0.75,-19.95}{naranja!70/publicación, grisr/registros y contabilidad,
  rojo!45/bus del sistema}{}

\end{tikzpicture}
\end{document}
